\documentclass[UTF8,11pt]{ctexart}
\usepackage[a4paper,margin=24mm]{geometry}
\usepackage{amsmath,amssymb,amsthm,mathtools}
\usepackage[all]{xy}
\usepackage{xurl,hyperref,enumitem}
\hypersetup{hidelinks,breaklinks=true}
\setlength{\emergencystretch}{3em}
\newtheorem*{theorem}{Theorem}
\newtheorem*{lemma}{Lemma}
\newtheorem*{proposition}{Proposition}
\newtheorem*{corollary}{Corollary}
\theoremstyle{definition}
\newtheorem*{definition}{Definition}
\newtheorem*{remark}{Remark}
\DeclareMathOperator{\Hom}{Hom}
\DeclareMathOperator{\Ext}{Ext}
\DeclareMathOperator{\Tor}{Tor}
\DeclareMathOperator{\Spec}{Spec}
\DeclareMathOperator{\supp}{supp}
\DeclareMathOperator{\im}{im}
\DeclareMathOperator{\id}{id}
\DeclareMathOperator{\Tr}{Tr}
\DeclareMathOperator{\rank}{rank}
\newcommand{\R}{\mathbb R}
\newcommand{\C}{\mathbb C}
\newcommand{\Z}{\mathbb Z}
\newcommand{\N}{\mathbb N}
\newcommand{\Q}{\mathbb Q}
\begin{document}
\section*{020\quad 流形紧对的完全相对Poincare对偶}
\subsection*{来源与收录范围}
Haynes Miller, \emph{Algebraic Topology I}, MIT 18.905, Fall 2016（基于 Sanath Devalapurkar 的课堂记录并经授课者修订）。主命题为 Theorem 37.1，印刷第102--104页；同时收录其明确回引的 Theorem 32.1 证明中第(3)、(5)步，以及 Lemma 32.4 的正文，印刷第89--90页。全部合为一个完全相对对偶问题。获取日期 2026-09-28。
\par\url{https://ocw.mit.edu/courses/18-905-algebraic-topology-i-fall-2016/2dad3dd8354992d877fbeca4049817ee_MIT18_905F16_lecture_notes.pdf}
\par\textbf{据原文整理的记号说明，非逐字引文：}$\check H$ 是由邻域对的奇异上同调取直极限得到的 \v Cech 上同调；$H_q(M|A)=H_q(M,M-A)$。$o_M$ 是定向局部系数系统。沿紧集 $K$ 的 $R$-定向，对应于限制到每个 $x\in K$ 都是局部同调生成元的基本类 $[M]_K\in H_n(M,M-K;R)$。定向类的存在（Theorem 32.1）、相对帽积的自然性及其与长正合列、Mayer--Vietoris 序列的兼容性（36.1--36.2）列作先修。另使用作者 Lemma 37.2 的紧集递减极限下 \v Cech 上同调连续性。
\subsection*{作者命题与人类证明}

\begin{theorem}[37.1, Fully relative Poincar\'e duality]
Let $M$ be an $n$-manifold and $K\supseteq L$ a pair of compact subsets. Assume given an $R$-orientation along $K$, with corresponding fundamental class $[M]_K$. With $p+q=n$, the map
\[
\cap[M]_K:\check H^p(K,L;R)\longrightarrow H_q(M-L,M-K;R)
\]
is an isomorphism.
\end{theorem}
We have seen that these isomorphisms are compatible; they form the rungs of the commuting ladder
\[
\xymatrix@C=8pt{
\cdots\ar[r]&\check H^{p-1}(L)\ar[r]\ar[d]^{\cap[M]_L}&\check H^p(K,L)\ar[r]\ar[d]^{\cap[M]_K}&\check H^p(K)\ar[r]\ar[d]^{\cap[M]_K}&\check H^p(L)\ar[r]\ar[d]^{\cap[M]_L}&\cdots\\
\cdots\ar[r]&H_{q+1}(M,M-L)\ar[r]&H_q(M-L,M-K)\ar[r]&H_q(M,M-K)\ar[r]&H_q(M,M-L)\ar[r]&\cdots
}
\]
To prove this theorem, we will follow the same five-step process we used to prove the Orientation Theorem 32.1. We have already prepared the Mayer--Vietoris ladder for this purpose.
\begin{proof}[Proof of Theorem 37.1]
By the top ladder and the five-lemma, we may assume $L=\varnothing$; so we want to prove that
\[
\cap[M]_K:\check H^p(K;R)\longrightarrow H_q(M,M-K;R)
\]
is an isomorphism.

(1) $M=\R^n$, $K$ a compact convex set. We claim that
\[
\check H^*(K)\xrightarrow{\cong}H^*(K).
\]
For any $\epsilon>0$, let $U_\epsilon$ denote the $\epsilon$-neighborhood of $K$,
\[
U_\epsilon=\bigcup_{x\in K}B_\epsilon(x).
\]
For any $y\in U_\epsilon$ there is a closest point in $K$, since the distance function to $y$ is continuous and bounded below on the compact set $K$ and so achieves its infimum. If $x',x''\in K$ are the same distance from $y$, then the midpoint of the segment joining $x'$ and $x''$ is closer, but lies in $K$ since $K$ is convex. So there is a unique closest point, $f(y)$. We let the listener check that $f:U_\epsilon\to K$ is continuous. It is also clear that if $i:K\to U_\epsilon$ is the inclusion then $i\circ f$ is homotopic to the identity on $Y$, by an affine homotopy.

Now let $D^n$ be a disk centered at the origin and containing the compact set $K$, and consider the commutative diagram
\[
\xymatrix{
H^p(K)\ar[r]^{\cap[\R^n]_K}&H_q(\R^n,\R^n-K)\\
H^p(D^n)\ar[u]^{\cong}\ar[r]\ar[d]_{\cong}&H_q(\R^n,\R^n-D^n)\ar[u]^{\cong}\ar[d]_{\cong}\\
H^p(*)\ar[r]&H_q(\R^n,\R^n-*).
}
\]
The groups are zero unless $p=0$, $q=n$. By naturality of the cap product, the bottom map is given by $1\mapsto1\cap[\R^n]_*$, and this is $[\R^n]_*$ since capping with $1$ is the identity, and this fundamental class is a generator of $H_n(\R^n,\R^n-*)$.

(2) $K$ a finite union of compact convex subsets of $\R^n$. This follows by induction and the five lemma applied to the Mayer--Vietoris ladder 36.2.

(3) $K$ is any compact subset of $\R^n$. This follows as before by a limit argument, using Lemmas 32.4 and 37.2.

(4) $M$ arbitrary, $K$ is a finite union of compact Euclidean subsets of $M$. This follows from (3) and Theorem 36.2.

(5) $M$ arbitrary, $K$ an arbitrary compact subset. This follows just as in the proof of Theorem 32.1.
\end{proof}
\paragraph{Referenced source material: Lemma 32.4}
\begin{lemma}[32.4]
Let $A_1\supseteq A_2\supseteq\cdots$ be a decreasing sequence of compact subsets of a space $X$, with intersection $A$. Then
\[
\varinjlim_i H_q(X,X-A_i)\xrightarrow{\cong}H_q(X,X-A).
\]
\end{lemma}
\begin{proof}
Let $\sigma:\Delta^q\to X$ be any $q$-simplex in $X-A$. The subsets $X-A_i$ form an open cover of $\im(\sigma)$, so by compactness it lies in some single $X-A_i$. This shows that
\[
\varinjlim_i S_q(X-A_i)\xrightarrow{\cong}S_q(X-A).
\]
Thus
\[
\varinjlim_i S_q(X|A_i)\xrightarrow{\cong}S_q(X|A_i)
\]
by exactness of direct limit, and the claim then follows for the same reason.
\end{proof}
\paragraph{Referenced source material: Theorem 32.1, steps (3) and (5)}
\noindent\textit{The source calls a subset of $M$ ``Euclidean'' if it lies inside some open set homeomorphic to $\R^n$.}

(3) For each $j\geq1$, let $C_j$ be a finite subset of $A$ such that
\[
A\subseteq\bigcup_{x\in C_j}B_{1/j}(x).
\]
Since any intersection of convex sets is either empty or convex,
\[
A_k=\bigcap_{j=1}^k\ \bigcup_{x\in C_j}B_{1/j}(x)
\]
is a union of finitely many convex sets, and since $A$ is closed it is the intersection of this decreasing family. So the result follows from (1), (2), and Proposition 32.3.

(5) Cover $A$ by finitely many open subsets that embed in Euclidean opens as open disks with compact closures. Their closures then form a finite cover by closed Euclidean disks $D_i$ in Euclidean opens $U_i$. For each $i$, excise the closed subset $M-U_i$ to see that
\begin{align*}
H_q(M,M-A\cap D_i)&\cong H_q(U_i,U_i-A\cap D_i)\\
&\cong H_q(\R^n,\R^n-A\cap D_i).
\end{align*}
By (4), the theorem holds for each of these. Each intersection $(A\cap D_i)\cap(A\cap D_j)$ is again a compact Euclidean subset, so the result holds for them by excision as well. The result then follows by (1).

\subsection*{整理说明（非作者正文）}
本条不只截取紧致定向流形的 Corollary 37.5；实际收录作者较强的完全相对命题及其全部五步证明，并把其回引的紧集近似与有限坐标覆盖段落取入，不能把这些段落另当两道题。先修只包括已明示的定向基本类、相对帽积、正合列及直极限性质；核心工作是把局部帽积同构经 Mayer--Vietoris、极限与坐标覆盖推广到一般紧对。

作者将最近点投影的连续性留给读者，本条保留这一提示并把它列作欧氏凸集的标准事实，而非宣称另有证明。PDF 文本中 Lemma 32.4 最后一式右端仍写 $A_i$；Theorem 37.1 写有 ``on $Y$''；所回引第(5)步末尾写 ``by (1)''，这里均未静默修订。原书的 ``convex'' 约定非空。印刷第104页的核心交换图已查看页面图像并重排；其他所取页的图像接口未成功，相关文字按取得的 PDF 文本转录。没有本地 PDF 原件。原书回引段落中(1)--(5)指 Theorem 32.1 的五步，不是本文件自动编号。
\subsection*{可修改的简短标注（非作者正文）}
$c$：流形紧子空间对的同调与上同调对偶。

$a$：局部定向系统；基本类；相对帽积；Cech 上同调；Mayer--Vietoris 序列；五引理；紧集近似；直极限；坐标覆盖与切除。

\texttt{cross\_theory\_status = yes}。通过帽积把上同调类变成低维同调类；局部欧氏几何中的基本类同构，经正合列、紧集极限和坐标拼接推广到任意紧对。位置：Theorem 37.1 的五步证明及其回引的 32.1(3)(5)。

全部标注均为非穷尽、可修改初稿；未标注不表示未使用。
\subsection*{许可与修改记录}
原文来自 MIT OpenCourseWare，作者与课程如上。按来源网站标示的 Creative Commons Attribution--NonCommercial--ShareAlike 4.0 许可复用；本转录及整理沿用同一许可。2026-09-28：助手转录、加入来源定位与简短标注；不暗示作者或 MIT 认可本次整理。许可：\url{https://creativecommons.org/licenses/by-nc-sa/4.0/}。
\end{document}
